\documentclass[../DoAn.tex]{subfiles}
\begin{document}

\begin{center}
    \Large{\textbf{TÓM TẮT NỘI DUNG ĐỒ ÁN}}\\
\end{center}
\vspace{1cm}

Trí tuệ nhân tạo lượng tử hóa là một hướng tiếp cận quan trọng để đưa các mô hình suy luận xuống hệ thống nhúng tài nguyên hạn chế. Trên các bộ xử lý RISC-V nhỏ, những phép toán INT8 như nhân--cộng tích lũy, tích vô hướng, kích hoạt và giới hạn miền giá trị thường phải được biểu diễn bằng nhiều lệnh vô hướng, làm tăng số lệnh thực thi và chi phí điều khiển vòng lặp. Các hướng tối ưu như cải thiện phần mềm, dùng SIMD/DSP, khai thác RISC-V Vector (RVV) hoặc bổ sung bộ tăng tốc chuyên dụng đều có thể đem lại lợi ích, nhưng cũng đặt ra đánh đổi về tài nguyên phần cứng, tính linh hoạt và khả năng tích hợp. Từ đó, đồ án lựa chọn hướng tiếp cận ở mức tập lệnh, nhằm giảm số lệnh và chi phí điều khiển ngay trong luồng thực thi của bộ xử lý, thay vì chuyển toàn bộ bài toán sang một bộ tăng tốc tách rời.

Trên cơ sở đó, đồ án đặc tả một nhóm thao tác phục vụ các kernel AI lượng tử hóa, gồm nhân--cộng tích lũy, tích vô hướng đóng gói, kích hoạt, giới hạn giá trị, tải dữ liệu hậu tăng và vòng lặp phần cứng. Các thao tác này được mô phỏng trên gem5 để kiểm tra ngữ nghĩa và đánh giá xu hướng hiệu năng. Song song với mô phỏng, đồ án xây dựng một hệ thống SoC tối giản sử dụng core CV32E40P trên kit FPGA DE2i-150, cấu hình phần mở rộng lệnh tương ứng trên core và xây dựng firmware để tạo các kịch bản đánh giá, bao gồm cả mô hình MNIST FC INT8 với TensorFlow Lite Micro. Kết quả cho thấy các kernel tối ưu trong mô phỏng có thể đạt mức tăng tốc hơn 9 lần ở cấu hình phù hợp, trong khi triển khai thực tế trên kit đạt khoảng 12--13 lần ở các bài đo chính so với đường tham chiếu phần mềm. Đóng góp chính của đồ án không chỉ nằm ở việc thiết kế tập lệnh RISC-V mở rộng cho suy luận AI lượng tử hóa, mà còn ở việc thiết lập một chuỗi kiểm chứng liền mạch từ mô phỏng đến triển khai thực tế. Qua đó, đồ án cho thấy tiềm năng của hướng tiếp cận này trên các nền tảng nhúng có tài nguyên hạn chế.

\begin{flushright}
Sinh viên thực hiện\\
\begin{tabular}{@{}c@{}}
\textit{(Ký và ghi rõ họ tên)}
\end{tabular}
\end{flushright}

\end{document}
